\begin{figure}
\centering
\resizebox{1\columnwidth}{!}{
\fbox{
\begin{minipage}[c][2.3in]{\linewidth}
% Use 'normalsize' font size and adjust line spacing
\normalsize
\begin{align*}
    & x = \texttt{'Le patient avait présenté une altération progressive de l'état général,}\\
    & \qquad \texttt{une fièvre et des sueurs nocturnes.'}\\
    & p = \text{concat}(t, x)\\
   & \begin{aligned}
        o = \Phi(p, \theta_h) = & \texttt{'-"fièvre" } \\ 
                                & \texttt{-"sueurs nocturnes"'}
      \end{aligned}\\
    & \begin{aligned}
        r(o, x) = & \; [ \\ 
                         & \space\space(\texttt{Le}, 0 , 2, O), (\texttt{patient}, 4 , 11, O), \dots, \\ 
                         & \space\space(\texttt{fièvre}, 77, 83, B_{clin}), (\texttt{et}, 85, 87, O), \\
                         & \space\space(\texttt{sueurs}, 91, 97, B_{clin}), (\texttt{nocturnes}, 98, 107, I_{clin}), \dots \\
                         & ]
      \end{aligned}
\end{align*}
\end{minipage}}
}
\caption{Illustration of our method's prediction and structuring steps on an example. The $t$ template is illustrated in Figure \ref{fig:prompt}. See glossary for translated example}
\label{fig:formalization}
\end{figure}
